% Resumos do percurso didatico; nao descrevem estudantes ou avaliacao reais.
\undfsetcontent{resumo}{%
Este relatório final apresenta um portfólio didático que reúne todos os projetos
de um percurso acadêmico fictício composto por duas experiências. O objetivo
foi relacionar uma entrega técnica e uma entrega documental por meio das
evidências necessárias para explicar seus resultados e limites. O primeiro
projeto produziu um validador de solicitações isoladas de reserva de laboratório,
com programa em Python e seis casos sintéticos. O segundo produziu um guia
para registrar contexto, contribuição individual, decisões, resultados e
materiais recuperáveis. A organização utilizou registros parciais, inventário
de projetos e identificadores de evidências, seguidos de sínteses por projeto
e de uma análise integradora. A execução do validador confirmou duas aceitações
e quatro recusas segundo as regras implementadas. O guia resultou em um
artefato documental organizado em seis tópicos, sem avaliação com leitores.
A análise relacionou rastreabilidade, delimitação de escopo e explicação dos
resultados, distinguindo a contribuição individual ilustrada da colaboração
atribuída à equipe fictícia. Concluiu-se que o conjunto permite demonstrar
como vincular afirmações a materiais verificáveis, sem comprovar implantação,
impacto ou evolução pessoal medida. A quantidade de projetos é apenas
demonstrativa e não estabelece requisito para a modalidade.}
\undfsetcontent{palavraschave}{portfólio; projetos aplicados; evidências}
\undfsetcontent{abstract}{%
This final report presents an illustrative portfolio containing all projects
in a fictional academic trajectory comprising two experiences. Its objective
was to connect a technical deliverable and a documentary deliverable through
the evidence needed to explain their findings and limitations. The first
project produced a validator for individual laboratory reservation requests,
including a Python program and six synthetic cases. The second produced a
guide for recording context, individual contribution, decisions, findings and
recoverable materials. The organization used partial records, a project
inventory and evidence identifiers, followed by project syntheses and an
integrative analysis. Executing the validator confirmed two acceptances and
four rejections under the implemented rules. The guide yielded a documentary
artifact organized into six topics, without evaluation by readers. The
analysis connected traceability, scope definition and explanation of findings,
distinguishing the illustrated individual contribution from collaboration
assigned to a fictional team. The complete set demonstrates how claims can
be linked to verifiable materials, without proving deployment, impact or
measured personal development. The number of projects serves only as an
example and does not establish a requirement for this academic format.}
\undfsetcontent{keywords}{portfolio; applied projects; evidence}
